% ---------------------------------------------------------------------------------------------------------------
% Section 4.3 -- Automatic Prompt Generator for SAM (replaces the previous Section 4.3).
% Requires: amsmath, graphicx, booktabs, adjustbox. \ResultsDir must point to the results/ directory.
% Tables are \input from \ResultsDir/tables (generated by prompter_bench/make_report.py), figures from
% \ResultsDir/figures (prompter_bench/figures.py; single panels in figures/panels/).
% ---------------------------------------------------------------------------------------------------------------

\section{Automatic Prompt Generator for SAM}
\label{sec:auto_prompt_results}

Phase~1 showed that the adapted SAM segments TCs and ARs well when it receives ground-truth prompts. An operational
system must obtain its prompts automatically, from a \emph{prompter}. This section asks whether such a pipeline
(prompter $\rightarrow$ prompts $\rightarrow$ SAM) outperforms the prompter alone, and which prompter and prompt format
work best. It is organised along five research questions:
\begin{description}
  \item[RQ1] How well can the frozen Phase-1 SAM segment TCs and ARs if the prompts are perfect, and how sensitive is it
  to realistic prompt errors? (Section~\ref{sec:prompt_analysis})
  \item[RQ2] Does SAM prompted by CG-Net improve on CG-Net, as reported in the earlier version of Table~4.13?
  (Section~\ref{sec:cgnet_prompter})
  \item[RQ3] Which prompter produces the best masks and prompts under identical conditions, and at what cost?
  (Section~\ref{sec:sam_input_generators})
  \item[RQ4] Does SAM add value to any prompter, and can this be achieved by training the prompter through SAM or by
  adapting the decoder? (Sections~\ref{sec:sam_input_generators} and~\ref{sec:decoder_adaptation_results})
  \item[RQ5] Where exactly do the automatic systems lose IoU? (Section~\ref{sec:error_analysis})
\end{description}
All experiments follow the protocol of Section~\ref{sec:prompting_protocol}: one frozen Phase-1 model (Infused Token,
MLP ratio 1.0), checkpoint selection on 40 held-out training images, three seeds per learned prompter, and evaluation
on the 61 test images.

% ===============================================================================================================
\subsection{Prompt Analysis}
\label{sec:prompt_analysis}

Before any prompter is trained, prompts built from the ground truth show what the frozen SAM can achieve and how it
reacts to the errors a real prompter makes.

\begin{figure}[htbp]
\centering
\includegraphics[width=0.8\linewidth]{\ResultsDir/figures/oracle_prompt_types.pdf}
\caption{SAM prompted with prompts built from the ground truth (test set, primary checkpoint).}
\label{fig:oracle_prompt_types}
\end{figure}

\paragraph{Prompt types.} With perfect prompts the tight box is the best sparse prompt (TC 0.723 / AR 0.625), points
are the weakest (0.527 / 0.577), and adding points to a box \emph{lowers} the IoU (0.583 / 0.604), as in
Table~4.12 (Figure~\ref{fig:oracle_prompt_types}). Dense prompts are better, because they carry the shape of the
object: mask logits alone reach AR~0.920, and the hybrid prompt (TC: box + mask, AR: mask) reaches TC~0.872 /
AR~0.920. These numbers are upper bounds, since a ground-truth mask prompt already contains the answer, but they show
that the adapted encoder and decoder are not the bottleneck of the automatic pipeline.

\begin{figure}[htbp]
\centering
\includegraphics[width=0.8\linewidth]{\ResultsDir/figures/oracle_mask_format.pdf}
\caption{How the decoder reads dense mask prompts built from the ground truth: $\{0,1\}$ maps versus logit maps of
scale $\pm a$, one prompt per class (union) versus one per object, with and without a box. Filled markers: primary
checkpoint (MLP ratio 1.0); open markers: MLP ratio 0.5.}
\label{fig:oracle_mask_format}
\end{figure}

\paragraph{Format of dense prompts.} Whether the decoder uses a dense prompt for TCs depends on the checkpoint
(Figure~\ref{fig:oracle_mask_format}). On the primary checkpoint both $\{0,1\}$ and logit masks work (TC 0.85,
AR~0.90); on the MLP-0.5 checkpoint a $\{0,1\}$ mask yields no TC at all (0.000) and even $\pm20$ logits reach only
TC~0.46, while ARs are unaffected. The prompt encoder reduces the $256\times256$ mask to $64\times64$, where a cyclone
covers only a handful of cells, and whether the decoder learned to use such weak dense evidence is a by-product of its
Phase-1 training. A box together with the mask is robust on both checkpoints (TC 0.86--0.87); this motivates the
hybrid conversion. One prompt per class works as well as one per object, so a single decoder call per class and image
suffices for dense prompts.

\begin{figure}[htbp]
\centering
\includegraphics[width=\linewidth]{\ResultsDir/figures/oracle_prompt_errors.pdf}
\caption{SAM with deliberately corrupted ground-truth prompts (primary checkpoint): (a)~box size, (b)~thickness of a
dense mask prompt, (c)~missed and false objects.}
\label{fig:oracle_prompt_errors}
\end{figure}

\begin{table}[htbp]
\centering
\small
\begin{adjustbox}{max width=\linewidth}
\begin{tabular}{lcccc}
\toprule
Ground-truth prompt & MLP 1.0 & MLP 0.5 & prompt-robust, hybrid & prompt-robust, hybrid (seed 1) \\
\midrule
tight box & 0.72 / 0.62 & 0.73 / 0.65 & 0.72 / 0.46 & 0.70 / 0.39 \\
box enlarged 20\% & 0.46 / 0.48 & 0.44 / 0.49 & 0.54 / 0.48 & 0.61 / 0.44 \\
box shrunk 20\% & 0.43 / 0.32 & 0.43 / 0.31 & 0.43 / 0.20 & 0.42 / 0.18 \\
box, 25\% of objects not prompted & 0.55 / 0.48 & 0.55 / 0.50 & 0.54 / 0.34 & 0.52 / 0.30 \\
box + 1 false box per class & 0.57 / 0.58 & 0.58 / 0.58 & 0.60 / 0.45 & 0.62 / 0.39 \\
box + 3 false boxes per class & 0.41 / 0.48 & 0.41 / 0.47 & 0.46 / 0.41 & 0.54 / 0.36 \\
mask logits $\pm$10 & 0.77 / 0.92 & 0.23 / 0.90 & 0.54 / 0.90 & 0.00 / 0.92 \\
mask eroded 5 px & 0.61 / 0.86 & 0.16 / 0.86 & 0.39 / 0.84 & 0.00 / 0.86 \\
mask dilated 5 px & 0.85 / 0.89 & 0.32 / 0.88 & 0.72 / 0.90 & 0.00 / 0.90 \\
box + mask logits $\pm$10 & 0.87 / 0.88 & 0.86 / 0.87 & 0.88 / 0.83 & 0.86 / 0.87 \\
\bottomrule
\end{tabular}
\end{adjustbox}
\caption{SAM prompted with prompts built from the ground truth, with controlled errors (TC IoU / AR IoU, 61 test images). Columns: the frozen Phase-1 decoders of both checkpoints and the decoders fine-tuned on generated or corrupted prompts (Section on decoder adaptation).}
\label{tab:oracle_prompts}
\end{table}


\paragraph{Sensitivity to prompt errors.} Figure~\ref{fig:oracle_prompt_errors} and Table~\ref{tab:oracle_prompts}
quantify how SAM reacts to the errors a prompter makes:
\begin{itemize}
  \item \emph{Box size.} The IoU peaks sharply at the tight box. Enlarging all boxes by 20\% costs 0.26 TC and 0.14 AR
  IoU; shrinking them by 20\% costs 0.30 and 0.31. The decoder takes the extent of the box literally.
  \item \emph{Missed and false objects.} An object without a prompt is never recovered (25\% of the objects without a
  prompt: TC $0.723\rightarrow0.548$, AR $0.625\rightarrow0.483$), and one random false box per class costs 0.149 TC and
  0.047 AR IoU: the decoder segments \emph{something} inside every box it receives.
  \item \emph{Dense prompts are forgiving.} Eroding the mask by 3~px costs only 0.07 TC / 0.02 AR, and dilating it by
  3--5~px even \emph{improves} TC ($0.774\rightarrow0.852$), because the $4\times$ downsampling in the prompt encoder
  otherwise weakens small objects.
\end{itemize}

\paragraph{Answer to RQ1.} With correct prompts the Phase-1 SAM is an excellent segmenter (hybrid prompts:
TC~0.87 / AR~0.92). It does not, however, correct its prompts: missed objects stay missed, false prompts become false
objects, and a box error of 10--20\% already costs 0.1--0.3 IoU.

% ===============================================================================================================
\subsection{SAM with CG-Net as Prompter}
\label{sec:cgnet_prompter}

\paragraph{Re-evaluation of Table~4.13.} Every configuration of the earlier Table~4.13 was re-run with the thesis
pipeline itself (the fine-tuned CG-Net, \texttt{PromptMaker}, \texttt{ClimateSAM.\allowbreak infer} and \texttt{forward}) on the
61 test images. Two image paths were compared: the path of \texttt{test\_\allowbreak prompt\_\allowbreak effect.\allowbreak py}, which calls
\texttt{ClimateSAM.\allowbreak set\_\allowbreak infer\_\allowbreak img()} and thereby skips the learned input adapter (it feeds the first three raw channels
to the encoder), and the training path \texttt{encode\_\allowbreak images()}, which applies it.

\input{\ResultsDir/tables/table_4_13_corrected}

\begin{figure}[htbp]
\centering
\includegraphics[width=0.8\linewidth]{\ResultsDir/figures/table_4_13_recheck.pdf}
\caption{SAM with CG-Net as prompter, all configurations of Table~4.13 (primary checkpoint). Filled: input adapter
applied; open: original script. Vertical line: CG-Net alone.}
\label{fig:table_4_13_recheck}
\end{figure}

The re-run of the original script with the checkpoint of the last recorded run reproduces that run exactly (box
prompts: TC 0.3404 / AR 0.3758 versus 0.3409 / 0.3758 in the stored result file), so the pipeline is faithfully
reproduced. The SAM rows of the earlier Table~4.13, however, are 0.07--0.12 IoU higher than every re-run and could not
be traced to any stored result; they most likely stem from an earlier state of the code.
Table~\ref{tab:cgnet_prompter_corrected} and Figure~\ref{fig:table_4_13_recheck} give the corrected values. The ranking
of the prompt types is confirmed (tight boxes $>$ boxes + points $>$ points; enlarging boxes hurts monotonically), but
\textbf{SAM prompted by CG-Net is never better than CG-Net itself}: the best configuration, tight boxes, reaches
TC~0.341 / AR~0.369 against 0.349 / 0.382 for CG-Net alone. The skipped input adapter is a genuine bug but changes the
IoU by only 0.00--0.03. The ``mask'' row depends on the checkpoint, as explained in Section~\ref{sec:prompt_analysis}.

\paragraph{CG-Net with a YOLO box head.} The thesis also attached a YOLO-style detection head to CG-Net that predicts
boxes directly. Its boxes are the weakest prompts of the study (Figure~\ref{fig:cgnet_yolo}; SAM: TC 0.19--0.22,
AR 0.22--0.29 for the three stored checkpoints and all objectness thresholds), far below boxes derived from CG-Net's
segmentation (0.34 / 0.37). Directly regressed boxes are too loose for thin, curved ARs, and the prompt analysis showed
that a box error of 20\% alone costs 0.14--0.26 IoU.

\begin{figure}[htbp]
\centering
\includegraphics[width=0.8\linewidth]{\ResultsDir/figures/cgnet_yolo.pdf}
\caption{SAM prompted by the boxes of the CG-Net YOLO head (three stored checkpoints A--C) as a function of the
objectness threshold. Grey line: CG-Net's own segmentation.}
\label{fig:cgnet_yolo}
\end{figure}

\paragraph{CG-Net trained under the protocol.} The fine-tuned CG-Net checkpoint used above was trained outside the
protocol of Section~\ref{sec:prompting_protocol}, on all 398 training images. For a fair comparison with the learned
SAM-feature prompters, CG-Net was re-trained on the 358 training images with its original Jaccard loss and the same
validation-based checkpoint selection, once starting from the official ClimateNet weights (two seeds, learning rate
$10^{-4}$) and once from random initialisation (three seeds, learning rate $10^{-3}$). Both are clearly better than the
stored checkpoint: TC~$0.350\pm0.005$ / AR~$0.399\pm0.005$ (mean FG IoU 0.375) from the official weights and
TC~$0.343\pm0.007$ / AR~$0.394\pm0.008$ (0.369) from scratch, against 0.349 / 0.382 (0.366) for the stored
checkpoint (Table~\ref{tab:prompters_infused_mlp1}). A carefully trained CG-Net on the four raw fields is therefore a
much stronger baseline than the earlier Table~4.13 suggested, and it is used as the reference prompter in
Section~\ref{sec:sam_input_generators}. SAM does not improve these masks either: with hybrid prompts the difference is
$-0.001$ mean FG IoU for both variants, with boxes $-0.007$ to $-0.009$.

\paragraph{Answer to RQ2.} The ranking of prompt types in Table~4.13 holds, its central claim does not: with correct
numbers, SAM prompted by CG-Net is slightly worse than CG-Net for every prompt type.

% ===============================================================================================================
\subsection{Generator from SAM Inputs}
\label{sec:sam_input_generators}

The prompters of this section read the frozen encoder's features instead of the raw climate fields: the linear probe
(logistic regression) on the first and on the last ViT block, the multi-scale fusion generator (MSF, Figure~3.9) and its
token-gated variants, and the mask-prompt generator (MPG, Section~\ref{sec:mpg}). The original multi-scale generators
predicted three channels, of which the loss supervised only TC and AR as independent sigmoids while the prediction was
the three-channel argmax, so an untrained background channel competed with TC and AR; here all prompters use two
sigmoid channels thresholded at zero. The original logistic-regression prompter read the \emph{first} ViT block; it is
evaluated as is and on the last block.

\begin{table}[htbp]
\centering
\small
\begin{tabular}{llccccc}
\toprule
Prompter & Output & TC IoU & AR IoU & BG IoU & Mean IoU & Mean FG IoU \\
\midrule
Ground truth (oracle) & Prompter mask (no SAM) & 1.000 & 1.000 & 1.000 & 1.000 & 1.000 \\
 & SAM + box & 0.723 & 0.625 & 0.971 & 0.773 & 0.674 \\
 & SAM + hybrid & 0.872 & 0.920 & 0.994 & 0.929 & 0.896 \\
\midrule
CG-Net (fine-tuned) & Prompter mask (no SAM) & 0.349 & 0.382 & 0.940 & 0.557 & 0.366 \\
 & SAM + box & 0.341 & 0.369 & 0.939 & 0.550 & 0.355 \\
 & SAM + hybrid & 0.349 & 0.384 & 0.941 & 0.558 & 0.367 \\
\bottomrule
\end{tabular}
\caption{Automatic prompting on the ClimateNet test set (61 images), frozen ClimateSAM (Infused Token, MLP ratio 1.0). Learned prompters: mean$\pm$std over 3 seeds, checkpoint selected on 40 held-out training images.}
\label{tab:prompters_infused_mlp1}
\end{table}


\begin{figure}[htbp]
\centering
\includegraphics[width=\linewidth]{\ResultsDir/figures/prompters_mask_vs_sam.pdf}
\caption{Prompter mask (grey) versus SAM prompted by it with boxes (orange), hybrid prompts (violet) and points
(yellow); learned static prompts (magenta) have no prompter. Error bars: standard deviation over three seeds.}
\label{fig:prompters_mask_vs_sam}
\end{figure}

\paragraph{Prompters.} Table~\ref{tab:prompters_infused_mlp1} and Figure~\ref{fig:prompters_mask_vs_sam} compare all
prompters. The mask-prompt generator gives the best prompter mask (TC $0.344\pm0.003$, AR $0.413\pm0.001$). The paired
bootstrap (Table~\ref{tab:bootstrap_selected}) locates its advantage over the fine-tuned CG-Net entirely in the ARs
(AR $+0.031$, 95\% CI $[+0.020,+0.043]$; TC $-0.005$, $[-0.030,+0.020]$); it is also better than MSF ($+0.011$ mean FG
IoU, $[+0.003,+0.019]$) and than the linear probe on the last block ($+0.036$, $[+0.025,+0.046]$). The token gate
\emph{lowers} MSF slightly but significantly ($-0.006$, $[-0.009,-0.003]$): its inputs are the fixed HQ tokens of the
frozen decoder, so the gate is a learned per-class channel weighting that carries no image-dependent information.
Label smoothing, which helped in Phase~1, does not help the prompter ($-0.004$, $[-0.007,-0.000]$). A linear probe on
the last ViT block (1.5\,k parameters) already reaches TC~0.301 / AR~0.385, close to the fine-tuned CG-Net, while the
same probe on the first block, the original configuration, reaches only 0.195 / 0.319: the late features of the adapted
encoder are almost linearly separable with respect to TC and AR, the early ones are not.

\paragraph{Comparison with the re-trained CG-Net.} Against CG-Net trained under the same protocol, the advantage of
MPG shrinks to the ARs (Table~\ref{tab:bootstrap_infused_mlp1}): AR $+0.018$ (95\% CI $[+0.009,+0.027]$) over CG-Net
from scratch and $+0.013$ ($[+0.003,+0.024]$) over CG-Net from the official weights, while TC is equal ($+0.001$ and
$-0.006$, both intervals spanning $\pm0.02$). The mean FG IoU difference ($+0.010$ and $+0.004$) is not significant.
The prompters built on the adapted SAM encoder thus match a well-trained dedicated network on the raw fields and are
slightly better for ARs, but they do not surpass it by a large margin.

\paragraph{Adding the raw fields to MPG.} Since CG-Net reaches nearly the same quality from only four raw fields, a
variant of MPG additionally receives these fields (average-pooled to $64\times64$ before the ConvNeXt blocks and resized
to $256\times256$ before the head, both branches initialised to zero). It is slightly \emph{worse} than MPG
(TC~$0.340\pm0.001$ / AR~$0.409\pm0.006$; $-0.004$ mean FG IoU, $[-0.008,-0.001]$). The adapted encoder already sees
all 16 input channels through the input adapter, so the four fields add no information that the SAM features lack.

\paragraph{Further token-gated variants.} Two further multi-scale generators of the thesis were run under the protocol.
The variant with context-guided (CG) blocks (1.31\,M parameters, 452~GFLOPs) trains slowly (1.5--3.5~hours per seed
on cached features, two seeds) and is clearly worse than the plain token gate: TC~$0.324\pm0.017$ /
AR~$0.355\pm0.022$ (mean FG IoU 0.340; $-0.023$, 95\% CI $[-0.032,-0.015]$), with its best validation epochs early
(epochs 2 and 14 of 60). Its AR mask is fragmented (20.5 AR objects per image), which is also the only case in which
SAM with hybrid prompts clearly improves on the prompter ($+0.017$, AR $+0.037$) -- the result, 0.356, is still below
the plain MPG mask. The variant with weights shared across the four block groups (0.39\,M parameters, 612~GFLOPs) is
slightly better than the token gate ($+0.006$, $[+0.003,+0.009]$; TC~$0.332\pm0.006$ / AR~$0.402\pm0.001$) but stays
below MPG ($-0.011$, $[-0.020,-0.004]$).

\paragraph{Box head on SAM features.} The thesis also designed a YOLO-style detection head on the image embedding,
gated by the decoder's HQ tokens (0.61\,M parameters), whose boxes prompt SAM directly. Trained on ground-truth boxes
(two seeds, objectness threshold 0.7 selected on the validation images), it is the weakest image-dependent prompter:
SAM with its boxes reaches TC~$0.107\pm0.017$ / AR~$0.303\pm0.025$ (mean FG IoU 0.205). It finds most objects (TC
recall 0.81--0.93), but SAM's output then contains 9--17 TC objects per image, of which only 17--29\% touch a
cyclone. As for the CG-Net YOLO
head, regressed boxes are too loose and too numerous; every false box becomes a false object
(Section~\ref{sec:prompt_analysis}).

\begin{table}[htbp]
\centering
\small
\begin{adjustbox}{max width=\linewidth}
\begin{tabular}{lccc}
\toprule
Comparison (B $-$ A) & $\Delta$ TC [95\% CI] & $\Delta$ AR [95\% CI] & $\Delta$ mean FG [95\% CI] \\
\midrule
MPG vs.\ CG-Net (fine-tuned), masks & $-0.005$ [$-0.030$, $+0.020$] & $\mathbf{+0.031}$ [$+0.020$, $+0.043$] & $+0.013$ [$-0.000$, $+0.027$] \\
MPG vs.\ MSF, masks & $+0.011$ [$-0.003$, $+0.025$] & $+0.010$ [$+0.005$, $+0.016$] & $\mathbf{+0.011}$ [$+0.003$, $+0.019$] \\
MPG vs.\ LogReg block 12, masks & $+0.044$ [$+0.025$, $+0.062$] & $+0.028$ [$+0.020$, $+0.036$] & $+0.036$ [$+0.025$, $+0.046$] \\
MSF + token gate vs.\ MSF, masks & $-0.009$ [$-0.015$, $-0.003$] & $-0.004$ [$-0.005$, $-0.002$] & $-0.006$ [$-0.009$, $-0.003$] \\
MPG: SAM + box vs.\ its mask & $-0.002$ [$-0.009$, $+0.006$] & $-0.022$ [$-0.028$, $-0.016$] & $-0.012$ [$-0.017$, $-0.007$] \\
MPG: SAM + hybrid vs.\ its mask & $-0.002$ [$-0.007$, $+0.003$] & $-0.015$ [$-0.024$, $-0.006$] & $-0.008$ [$-0.014$, $-0.003$] \\
CG-Net (ft.): SAM + box vs.\ its mask & $-0.008$ [$-0.017$, $+0.001$] & $-0.013$ [$-0.018$, $-0.006$] & $-0.010$ [$-0.016$, $-0.004$] \\
MPG end-to-end: SAM + hybrid vs.\ its mask & $-0.001$ [$-0.005$, $+0.003$] & $+0.024$ [$+0.021$, $+0.028$] & $+0.012$ [$+0.009$, $+0.015$] \\
MPG two-stage: SAM + hybrid vs.\ MPG mask & $-0.007$ [$-0.013$, $-0.000$] & $-0.002$ [$-0.007$, $+0.003$] & $-0.004$ [$-0.008$, $-0.000$] \\
\bottomrule
\end{tabular}
\end{adjustbox}
\caption{Selected paired-bootstrap comparisons on the test set (2000 resamples of the 61 images, seeds pooled). The
full list is given in Table~\ref{tab:bootstrap_infused_mlp1}.}
\label{tab:bootstrap_selected}
\end{table}

\begin{figure}[htbp]
\centering
\includegraphics[width=\linewidth]{\ResultsDir/figures/sam_minus_prompter.pdf}
\caption{Effect of SAM: IoU of SAM's output minus IoU of the prompter's own mask, per prompter and prompt type, with
paired-bootstrap 95\% intervals.}
\label{fig:sam_minus_prompter}
\end{figure}

\paragraph{Added value of SAM.} Figure~\ref{fig:sam_minus_prompter} answers the central question of Phase~2. Point
prompts always lose ($-0.03$ to $-0.08$ mean FG IoU). Tight boxes are roughly neutral for TCs ($-0.010$ to $+0.011$)
and lose 0.013--0.024 AR IoU for every prompter with a good AR mask; they are neutral or slightly positive only where
the prompter's AR mask is poor (official CG-Net, linear probe on block~1, end-to-end MPG). The hybrid prompt is the best
conversion: within about $\pm0.01$ of the prompter mask for every prompter, slightly \emph{positive} for prompters with a
poor mask (linear probe on block~1: $+0.007$; end-to-end MPG: $+0.012$) and slightly \emph{negative} for the strong ones
(MSF $-0.009$, MPG $-0.008$). The best SAM output of the study (two-stage MPG, hybrid prompts, 0.374 mean FG IoU) does
not exceed the best prompter mask (MPG, 0.379). The complete comparison of all seven conversions is given in
Table~\ref{tab:prompt_conversion_infused_mlp1}.

\paragraph{Training through SAM.} Training MPG end-to-end through the frozen decoder makes \emph{its own} mask worse
(0.356 versus 0.379) and SAM's output better than that mask ($+0.012$): the prompter learns to produce logits that SAM
turns into good masks. The resulting SAM output (0.368) nevertheless stays below the plain MPG mask
($-0.011$, $[-0.016,-0.005]$), and the end-to-end mask is strongly fragmented (Table~\ref{tab:object_level_infused_mlp1}),
because a loss on SAM's output does not penalise speckle that SAM ignores. The two-stage variant keeps the MPG mask
quality (0.374) and gives the best SAM output of the study (0.374), again not above the MPG mask ($-0.004$,
$[-0.008,-0.000]$). Its selected validation epochs are 2--8 of 30: once the prompter has been trained on its own maps,
training through SAM barely changes it.

\begin{figure}[htbp]
\centering
\includegraphics[width=0.75\linewidth]{\ResultsDir/figures/efficiency.pdf}
\caption{Mean FG IoU of the prompter mask versus forward cost per image (log scale).}
\label{fig:efficiency}
\end{figure}

\begin{table}[htbp]
\centering
\small
\begin{adjustbox}{max width=\linewidth}
\begin{tabular}{lccc}
\toprule
Model & Parameters & GFLOPs / image & Latency (ms, A100, batch 1) \\
\midrule
CG-Net & 494,232 & 22.93 & 8.89 \\
Logistic regression & 1,538 & 0.01 & 0.05 \\
Multi-scale fusion & 1,406,474 & 702.36 & 36.13 \\
Multi-scale fusion + token gate & 1,415,180 & 702.89 & 38.12 \\
Mask-prompt generator & 671,138 & 5.74 & 0.89 \\
(frozen ClimateSAM encoder, for reference) & 95,987,456 & 973.25 & 122.39 \\
(one SAM decoder call, one box, for reference) & 5,269,508 & 16.80 & 3.39 \\
\bottomrule
\end{tabular}
\end{adjustbox}
\caption{Cost of the prompters on top of the frozen ClimateSAM encoder (forward pass, one image; FLOPs counted with torch.utils.flop\_counter, latency = median of 20 runs). The encoder and one decoder call are listed for reference.}
\label{tab:speed}
\end{table}


\paragraph{Cost.} MPG is both the most accurate and one of the cheapest prompters (Figure~\ref{fig:efficiency},
Table~\ref{tab:speed}): 5.7~GFLOPs and 0.9~ms per image, i.e.\ $122\times$ fewer FLOPs and $40\times$ lower latency than
MSF and $4\times$ fewer FLOPs than CG-Net. Relative to the frozen encoder (973~GFLOPs) the prompter adds 0.6\%. MSF
spends almost all of its computation on $3\times3$ convolutions at $1024\times1024$ over features that were upsampled
by nearest-neighbour interpolation from $64\times64$, i.e.\ on a resolution that contains no additional information.
With cached features a complete MPG training takes about 4~minutes and converges within 10--15 epochs
(Figure~\ref{fig:training_curves}); an MSF training takes 46--83 minutes.

\begin{figure}[htbp]
\centering
\includegraphics[width=\linewidth]{\ResultsDir/figures/training_curves.pdf}
\caption{Validation mean FG IoU of the prompter mask and training loss (mean $\pm$ std over three seeds). Loss values
are not comparable across architectures, because the MSF models add deep-supervision terms.}
\label{fig:training_curves}
\end{figure}

\paragraph{Answer to RQ3 and RQ4.} The best single prompters are the small mask-prompt generator and a CG-Net trained
under the same protocol; MPG is better for ARs and far cheaper than the other SAM-feature prompters. SAM prompted by
any prompter is at best as good as that prompter, and training the prompter through SAM does not change this.

% ===============================================================================================================
\subsection{Error Analysis}
\label{sec:error_analysis}

\begin{table}[htbp]
\centering
\small
\begin{tabular}{lcccccc}
\toprule
Prompter & TC recall & TC precision & AR recall & AR precision & TC obj./img & AR obj./img \\
\midrule
Ground truth (oracle) & 1.000 & 1.000 & 1.000 & 1.000 & 2.9 & 6.8 \\
CG-Net (fine-tuned) & 0.682 & 0.513 & 0.908 & 0.567 & 3.9 & 11.9 \\
\bottomrule
\end{tabular}
\caption{Object-level detection of the prompter masks: a ground-truth object is found if any predicted pixel overlaps it; a predicted object is correct if it overlaps any ground-truth object. The ground truth has 2.6 TC and 7.4 AR objects per image (Table 2.3).}
\label{tab:object_level_infused_mlp1}
\end{table}


\begin{figure}[htbp]
\centering
\includegraphics[width=\linewidth]{\ResultsDir/figures/object_level.pdf}
\caption{Object-level recall and precision of the prompter masks (mean over seeds; bars: standard deviation).}
\label{fig:object_level}
\end{figure}

\paragraph{Detection.} All prompters find more than 90\% of the ARs but only about two thirds of the TCs, and 40--60\%
of their predicted objects touch no ground-truth object (Table~\ref{tab:object_level_infused_mlp1},
Figure~\ref{fig:object_level}). For ARs many of these are small fragments: the prompters predict 12--18 AR objects per
image, against 6.8 in the ground truth.

\begin{figure}[htbp]
\centering
\includegraphics[width=\linewidth]{\ResultsDir/figures/error_decomposition.pdf}
\caption{IoU of the prompter masks as predicted, with perfect detection (false objects removed, missed objects added),
and with the true shape of every detected object (false objects kept).}
\label{fig:error_decomposition}
\end{figure}

\paragraph{Detection versus delineation.} The error decomposition (Figure~\ref{fig:error_decomposition},
Table~\ref{tab:decomposition_infused_mlp1}) separates the two error sources.
\begin{itemize}
  \item \emph{ARs.} For MPG, perfect detection raises the AR IoU only from 0.413 to 0.490, but giving every
  \emph{detected} river its true shape raises it to 0.817. The rivers are found; what is wrong is their extent, i.e.\
  their length, their width and where a river ends.
  \item \emph{TCs.} Perfect detection ($0.345\rightarrow0.585$) and perfect shapes of the detected cyclones
  ($\rightarrow0.636$) matter about equally: a third of the cyclones is missed, half of the predicted blobs are false,
  and the small objects are also delineated imprecisely.
\end{itemize}
These are exactly the errors that SAM cannot repair (Section~\ref{sec:prompt_analysis}): missed objects stay missed,
false prompts become false objects, and the extent of an AR is taken from the prompt. This is the mechanism behind the
result of Section~\ref{sec:sam_input_generators}.

\paragraph{Answer to RQ5.} ARs are detected ($>90\%$) but their extent is wrong; TCs are missed (one third) and falsely
detected (about half of the predicted blobs).

% ===============================================================================================================
\subsection{Adapting the Decoder to Generated Prompts}
\label{sec:decoder_adaptation_results}

\begin{figure}[htbp]
\centering
\includegraphics[width=\linewidth]{\ResultsDir/figures/decoder_adaptation.pdf}
\caption{(a)~Test mean FG IoU of the prompter mask, SAM with the Phase-1 decoder and SAM with the adapted decoder;
(b)~validation IoU during fine-tuning (epoch~0: Phase-1 decoder; the CG-Net curve is high because CG-Net was trained on
these images).}
\label{fig:decoder_adaptation}
\end{figure}

\begin{table}[htbp]
\centering
\small
\begin{adjustbox}{max width=\linewidth}
\begin{tabular}{llccccc}
\toprule
Prompter, prompts, training prompts & Output & TC IoU & AR IoU & Mean FG IoU & TC recall / prec. & AR recall / prec. \\
\midrule
CG-Net (fine-tuned), bbox, in-sample & prompter mask & 0.349 & 0.382 & 0.366 &  &  \\
 & SAM, Phase-1 decoder & 0.341 & 0.369 & 0.355 & 0.67 / 0.56 & 0.90 / 0.53 \\
 & SAM, adapted decoder & 0.345 & 0.384 & 0.364 & 0.64 / 0.58 & 0.88 / 0.58 \\
\midrule
Mask-prompt generator, bbox, out-of-fold & prompter mask & 0.346 & 0.414 & 0.380 &  &  \\
 & SAM, Phase-1 decoder & 0.344 & 0.392 & 0.368 & 0.65 / 0.59 & 0.95 / 0.50 \\
 & SAM, adapted decoder & 0.342 & 0.404 & 0.373 & 0.65 / 0.57 & 0.93 / 0.62 \\
\midrule
Mask-prompt generator, bbox, in-sample & prompter mask & 0.346 & 0.414 & 0.380 &  &  \\
 & SAM, Phase-1 decoder & 0.344 & 0.392 & 0.368 & 0.65 / 0.59 & 0.95 / 0.50 \\
 & SAM, adapted decoder & 0.343 & 0.404 & 0.374 & 0.65 / 0.59 & 0.95 / 0.55 \\
\midrule
Mask-prompt generator, hybrid, out-of-fold & prompter mask & 0.346 & 0.414 & 0.380 &  &  \\
 & SAM, Phase-1 decoder & 0.344 & 0.397 & 0.371 & 0.65 / 0.59 & 0.91 / 0.65 \\
 & SAM, adapted decoder & 0.347 & 0.405 & 0.376 & 0.65 / 0.60 & 0.94 / 0.64 \\
\bottomrule
\end{tabular}
\end{adjustbox}
\caption{Adapting the HQ decoder parts to generated prompts (test set). Phase-1 decoder = the frozen decoder of all other experiments; adapted = best validation epoch of 30.}
\label{tab:decoder_adaptation_infused_mlp1}
\end{table}


\paragraph{Adaptation to a prompter.} Fine-tuning the decoder on generated prompts (Section~\ref{sec:decoder_adaptation_method})
teaches it something useful: with out-of-fold MPG boxes it rejects part of the false AR prompts (AR object precision
$0.50\rightarrow0.62$) and it recovers most of what SAM loses relative to the prompter (Figure~\ref{fig:decoder_adaptation},
Table~\ref{tab:decoder_adaptation_infused_mlp1}). In no configuration, however, does SAM end up above the prompter mask
(best: $-0.004$ mean FG IoU with out-of-fold hybrid prompts); on the second checkpoint the best configuration ties the
prompter ($+0.001$).

\paragraph{Prompt-robust decoder.} The adaptation above changes only which prompts the decoder sees. The prompt-robust
decoder (Section~\ref{sec:decoder_adaptation_method}) is additionally trained on deliberately corrupted ground-truth
prompts whose targets are the complete objects, and empty masks for false boxes. Table~\ref{tab:robust_decoder_infused_mlp1}
and Figure~\ref{fig:robust_decoder} report each decoder on clean and corrupted ground-truth prompts (as in
Table~\ref{tab:oracle_prompts}) and with the prompts of MPG and of the fine-tuned CG-Net.

\begin{table}[htbp]
\centering
\small
\begin{adjustbox}{max width=\linewidth}
\begin{tabular}{lcccccccc}
\toprule
Decoder & GT box & GT box +20\% & GT box + 1 false & GT box + mask & MPG: box & MPG: hybrid & CG-Net: box & CG-Net: hybrid \\
\midrule
Phase-1 decoder & 0.674 & 0.474 & 0.576 & 0.877 & 0.367 & 0.370 & 0.355 & 0.367 \\
prompt-robust, hybrid & 0.592 & 0.511 & 0.524 & 0.856 & 0.332 & 0.375 & 0.320 & 0.367 \\
prompt-robust, hybrid (seed 1) & 0.548 & 0.527 & 0.503 & 0.867 & 0.317 & 0.376 & 0.309 & 0.369 \\
\bottomrule
\end{tabular}
\end{adjustbox}
\caption{Mean FG IoU on the test set of SAM with replaced decoders. GT: prompts built from the ground truth (clean or corrupted, as in the oracle study); MPG: SAM prompted by the mask-prompt generator (mean of three seeds; its own mask reaches 0.379); CG-Net: fine-tuned CG-Net (own mask 0.366).}
\label{tab:robust_decoder_infused_mlp1}
\end{table}


\begin{figure}[htbp]
\centering
\includegraphics[width=\linewidth]{\ResultsDir/figures/robust_decoder.pdf}
\caption{Prompt-robust decoders (mean of two seeds) versus the Phase-1 decoder. (a,~b)~Ground-truth prompts with
controlled errors; (c)~prompts of the mask-prompt generator (mean of three seeds); dashed: the generator's own mask.}
\label{fig:robust_decoder}
\end{figure}

\begin{itemize}
  \item \emph{The decoder does become more robust to box errors.} Trained in box mode, it loses less on enlarged boxes
  (ground-truth boxes enlarged by 20\%: mean FG IoU $0.474\rightarrow0.559$--$0.564$) and on false boxes
  ($0.576\rightarrow0.608$--$0.609$), at the price of tight boxes ($0.674\rightarrow0.630$--$0.636$): it no longer takes
  the box extent literally.
  \item \emph{This robustness hardly transfers to real prompters.} SAM with MPG boxes improves by only $+0.005$
  ($[+0.002,+0.008]$, both seeds), with CG-Net boxes by $+0.008$--$0.010$. Trained in hybrid mode, the decoder improves
  SAM with MPG hybrid prompts by $+0.005$--$0.006$ ($[+0.004,+0.007]$), to 0.375--0.376, which is still below the MPG mask
  itself ($-0.003$ to $-0.004$, $[-0.006,-0.000]$). This is the same gain as the adaptation to out-of-fold MPG prompts
  alone.
  \item \emph{The decoder specialises to its prompt mode.} The hybrid-trained decoder is much worse with box-only prompts
  (MPG boxes: $-0.035$ to $-0.050$), the box-trained one with hybrid prompts ($-0.006$ to $-0.010$).
  \item \emph{Synthetic corruption alone does not help.} Without prompter prompts, the validation IoU with MPG prompts
  falls in every epoch, so checkpoint selection keeps the Phase-1 decoder (epoch~0). The decoder of the last epoch is
  worse with real prompts ($-0.010$ with MPG hybrid prompts, $-0.076$ with MPG boxes), while on ground-truth prompts of
  its own mode it is unchanged (box + mask: 0.881 versus 0.877; Table~\ref{tab:robust_decoder_infused_mlp1}). Hand-designed corruptions do not match the
  errors of a real prompter.
\end{itemize}
Training the decoder for robustness therefore works as intended on the synthetic errors but does not lift SAM above
its prompter: the real prompter errors that remain after training (missed cyclones and the extent of rivers,
Section~\ref{sec:error_analysis}) cannot be corrected from a prompt that lacks the information.

% ===============================================================================================================
\subsection{Evaluation-only Improvements}
\label{sec:posthoc}

Three standard evaluation-time techniques were tested on top of the trained prompters, with every free parameter
chosen on the 40 validation images (Table~\ref{tab:posthoc_infused_mlp1}): averaging the logits of the three seeds
(\emph{ensemble}), a per-class logit threshold instead of zero (\emph{calibrated}), and object-level post-processing
with a minimum blob size per class and a maximum latitude for TCs (\emph{post-processed}). Finally, SAM's own output was
fed back as the next prompt for up to three rounds, SAM's native interactive refinement loop.

\input{\ResultsDir/tables/posthoc_infused_mlp1}

\begin{itemize}
  \item \emph{Seed ensembles} help every non-linear prompter by 0.006--0.011 mean FG IoU, and end-to-end MPG, whose
  seeds differ most, by 0.015 (MPG: $0.379\rightarrow0.385$,
  95\% CI of the gain over seed~0 $[+0.001,+0.010]$; MSF: $0.368\rightarrow0.376$; CG-Net from scratch:
  $0.369\rightarrow0.380$). For the SAM-feature prompters the gain comes mostly from TCs ($+0.007$ to $+0.011$
  against $+0.004$ to $+0.007$ for ARs), for CG-Net from both classes (TC $+0.009$, AR $+0.014$). The linear probe gains nothing, as
  its three seeds converge to nearly the same classifier.
  \item \emph{Threshold calibration and post-processing} do not transfer from the validation to the test images. The
  thresholds selected on the 40 validation images are far from zero for TCs (logits between $-1.0$ and $+2.75$,
  depending on the prompter), but on the test set they change the IoU by only $-0.005$ to $+0.002$. With so few
  validation images the selection mostly fits noise; the unbiased threshold of the loss is the better choice.
  \item \emph{Iterative SAM refinement} lowers the IoU in every round (MPG: $0.384\rightarrow0.369\rightarrow0.366
  \rightarrow0.359$). Each round loses a little AR extent and some TCs and recovers nothing, which is the prompt
  analysis of Section~\ref{sec:prompt_analysis} applied repeatedly.
  \item \emph{Ensembles across prompter families.} Averaging the logits of the three MPG and the three CG-Net
  (from scratch) models gives the best mask of the study, TC~0.363 / AR~0.416 (mean FG IoU 0.389; 0.391 with calibrated
  thresholds). Relative to the MPG ensemble the gain lies in the TCs ($+0.010$, CI $[-0.005,+0.024]$), relative to the
  CG-Net ensemble in both classes ($+0.010$ mean FG IoU, $[+0.004,+0.016]$). Prompters on the SAM features and on the raw
  fields thus make partly different TC errors. The comparison with the three-model ensembles is not entirely fair (six
  models against three), and the gain over the MPG ensemble is not significant ($+0.004$, $[-0.005,+0.012]$).
\end{itemize}

% ===============================================================================================================
\subsection{Prompt-free SAM with Learned Static Prompts}
\label{sec:static_prompts}

The most minimal prompt-free variant needs no prompter: $K$ sparse prompt tokens per class (initialised from SAM's
positive-point embedding) are learned and fed to the frozen decoder for every image, so all image-specific information
must come from the decoder's attention to the image embedding. With $K=4$ this reaches TC~$0.237\pm0.001$ /
AR~$0.324\pm0.008$ (three seeds), with $K=16$ TC~0.248 / AR~0.329, i.e.\ 0.09--0.10 mean FG IoU below MPG. Object
recall is high (TC~0.87, AR~0.98), but the output breaks into thousands of fragments (object precision about 0.2):
a fixed query cannot decide \emph{which} humid, windy region is a TC or an AR. This decision needs an image-dependent
prompter.

% ===============================================================================================================
\subsection{Robustness to the Frozen Checkpoint}
\label{sec:second_checkpoint}

\begin{figure}[htbp]
\centering
\includegraphics[width=0.55\linewidth]{\ResultsDir/figures/second_checkpoint.pdf}
\caption{Mean FG IoU of every prompter and output on the MLP-1.0 ($x$) and MLP-0.5 ($y$) checkpoints. Diagonal: equal
performance.}
\label{fig:second_checkpoint}
\end{figure}

All experiments were repeated with the MLP-0.5 checkpoint (Table~\ref{tab:prompters_infused_mlp05}; MSF models with one
seed). The results lie close to the diagonal (Figure~\ref{fig:second_checkpoint}), so the conclusions do not depend on
the checkpoint. On this checkpoint MPG (0.370) ties MSF (0.370) and the smoothed MPG (0.371); all three stay ahead of
CG-Net (0.366; AR $+0.027$ to $+0.028$) and of the linear probe (0.342), and again no SAM output exceeds the best
prompter mask (best SAM output 0.373 versus 0.371, within noise).

% ===============================================================================================================
\subsection{Qualitative Examples}
\label{sec:qualitative}

\begin{figure}[htbp]
\centering
\includegraphics[width=\linewidth]{\ResultsDir/figures/maps_test_image_17.pdf}
\caption{Test image 17. Filled: ground truth (TC green, AR blue); outlines: prediction; black rectangles: box prompts.
(b--d)~prompter masks; (e,~f)~SAM prompted by MPG with boxes and with hybrid prompts; (g)~SAM prompted by CG-Net boxes;
(h)~SAM with learned static prompts.}
\label{fig:maps_17}
\end{figure}

Figure~\ref{fig:maps_17} illustrates the quantitative findings. All prompters find the long rivers of the southern
hemisphere and the large cyclone over the Gulf of Mexico; they differ in the extent of the rivers and in small false or
missed objects. SAM with boxes (e,~g) fills every box with a river-shaped mask, including the false AR over the central
North Pacific that the prompters hallucinated, and cuts rivers at the box borders; the hybrid prompt (f) follows the
prompter's shape closely. The learned static prompts (h) mark every moist band. A second example is shown in
Figure~\ref{fig:maps_45} (Appendix).

% ===============================================================================================================
\subsection{Summary}
\label{sec:auto_prompt_summary}

\begin{itemize}
  \item \textbf{RQ1.} With perfect prompts the Phase-1 SAM reaches TC~0.87 / AR~0.92 (hybrid prompts). It does not
  correct prompts: missed objects stay missed, false prompts become false objects, and 10--20\% box error costs
  0.1--0.3 IoU.
  \item \textbf{RQ2.} The SAM rows of the earlier Table~4.13 do not reproduce; SAM prompted by CG-Net is slightly worse
  than CG-Net for every prompt type. The ranking of the prompt types is confirmed.
  \item \textbf{RQ3.} The mask-prompt generator (0.67\,M parameters, 5.7~GFLOPs) is the best prompter: TC~0.344 /
  AR~0.413, AR $+0.031$ over the stored fine-tuned CG-Net and $+0.011$ mean FG IoU over the multi-scale fusion
  generator at $122\times$ fewer FLOPs. Against CG-Net re-trained under the same protocol it is better only for ARs
  ($+0.013$ to $+0.018$); the overall difference is not significant. An ensemble of both families gives the best mask
  (0.389).
  \item \textbf{RQ4.} SAM prompted by any automatic prompter is at best as good as that prompter. Neither training the
  prompter through SAM nor adapting the decoder to realistic prompts changes this.
  \item \textbf{RQ5.} ARs are found but their extent is wrong; TCs are missed and falsely detected. A prompt-conditioned
  decoder cannot repair either error.
\end{itemize}
